%% LyX 2.3.4.3 created this file.  For more info, see http://www.lyx.org/.
%% Do not edit unless you really know what you are doing.
\documentclass[oneside,english]{amsart}
\usepackage{amsbsy}
\usepackage{amstext}
\usepackage{amsthm}
\usepackage{amssymb}
\usepackage{fontspec}
\usepackage{geometry}
\geometry{verbose,tmargin=1cm,bmargin=1.5cm,lmargin=1cm,rmargin=1cm}
\setlength{\parskip}{\smallskipamount}
\setlength{\parindent}{0pt}
\usepackage{xcolor}
\definecolor{shadecolor}{rgb}{0.9375, 0.9375, 0.9375}
\usepackage{framed}
\usepackage{mathrsfs}
\usepackage{graphicx}
\usepackage{setspace}
\usepackage[unicode=true]
 {hyperref}

\makeatletter

%%%%%%%%%%%%%%%%%%%%%%%%%%%%%% LyX specific LaTeX commands.
\providecommand\textquotedblplain{%
  \bgroup\addfontfeatures{Mapping=}\char34\egroup}
%% Because html converters don't know tabularnewline
\providecommand{\tabularnewline}{\\}

%%%%%%%%%%%%%%%%%%%%%%%%%%%%%% Textclass specific LaTeX commands.
\numberwithin{equation}{section}
\numberwithin{figure}{section}
\theoremstyle{plain}
\newtheorem{thm}{\protect\theoremname}
\theoremstyle{plain}
\newtheorem{conjecture}[thm]{\protect\conjecturename}
\theoremstyle{definition}
\newtheorem*{defn*}{\protect\definitionname}

%%%%%%%%%%%%%%%%%%%%%%%%%%%%%% User specified LaTeX commands.
\usepackage{shuffle}

\makeatother

\usepackage{polyglossia}
\setdefaultlanguage[variant=american]{english}
\providecommand{\conjecturename}{Conjecture}
\providecommand{\definitionname}{Definition}
\providecommand{\theoremname}{Theorem}

\begin{document}
\begin{center} {\Large  Conjecture on the Number of Linearly Independent Relations of Ihara--Kaneko--Zagier Derivation Relation: An Euler-type Pentagonal Formula } \\[1.5em]
{\bf  Tamás Szűcs} \\[0.5em]{October 05, 2026 -- Szeged} \\[0.5em] {\tt szucs.tamas.prof@gmail.com} \end{center}

\section{Generalized Pentagonal Numbers and the $\omega(n)$ Function}

Let $P_{0},P_{1},P_{2},P_{3},P_{4},P_{5},P_{6},P_{7},\ldots=0,1,2,5,7,12,15,22,\ldots$
denote the generalized pentagonal numbers \href{https://oeis.org/A001318}{OEIS:A001318},
and let $\omega(k)$= \href{https://oeis.org/A010815}{A010815(k)},
defined by 
\[
\omega(k)=\begin{cases}
1 & {\text{if }}k=0;\\
(-1)^{m} & {\text{if }}k=\frac{3m^{2}\pm m}{2};\\
0 & {\text{otherwise.}}
\end{cases}
\]

\parbox[t]{1\columnwidth}{%
\begin{center}
Table 1. $P_{n}$ and $\omega(n)$
\par\end{center}
\begin{center}
{\small{}}%
\begin{tabular}{|c|c|c|c|c|c|c|c|c|c|c|c|c|c|c|c|c|c|c|c|c|c|c|c|c|}
\hline 
{\small{}$\vphantom{\dfrac{A}{B}}n$} & {\small{}0} & {\small{}1} & {\small{}2} & {\small{}3} & {\small{}4} & {\small{}5} & {\small{}6} & {\small{}7} & {\small{}8} & {\small{}9} & {\small{}10} & {\small{}11} & {\small{}12} & {\small{}13} & {\small{}14} & {\small{}15} & {\small{}16} & {\small{}17} & {\small{}18} & {\small{}19} & {\small{}20} & {\small{}21} & {\small{}22} & {\small{}23}\tabularnewline
\hline 
{\small{}$\vphantom{\dfrac{A}{B}}P_{n}$} & {\small{}0} & {\small{}1} & {\small{}2} & {\small{}5} & {\small{}7} & {\small{}12} & {\small{}15} & {\small{}22} & {\small{}26} & {\small{}35} & {\small{}40} & {\small{}51} & {\small{}57} & {\small{}70} & {\small{}77} & {\small{}92} & {\small{}100} & {\small{}117} & {\small{}126} & {\small{}145} & {\small{}155} & {\small{}176} & {\small{}187} & {\small{}210}\tabularnewline
\hline 
{\small{}$\vphantom{\dfrac{A}{B}}\omega(n)$} & {\small{}1} & {\small{}-1} & {\small{}-1} & {\small{}0} & {\small{}0} & {\small{}1} & {\small{}0} & {\small{}1} & {\small{}0} & {\small{}0} & {\small{}0} & {\small{}0} & {\small{}-1} & {\small{}0} & {\small{}0} & {\small{}-1} & {\small{}0} & {\small{}0} & {\small{}0} & {\small{}0} & {\small{}0} & {\small{}0} & {\small{}1} & {\small{}0}\tabularnewline
\hline 
\end{tabular}{\small\par}
\par\end{center}%
}

\section{The Conjecture and Some Equivalent Formulations }

For a given weight $n\geq3$ let $\mathscr{L}(n)$ denote the the
number of linearly independent relations of Ihara-Kaneko-Zagier derivation
relation. In \cite{doi:10.1142/S1793042117500245,TANAKA20092021}
the authors have tabulated the values of $\mathscr{L}(n)$ calculated
using the Risa/Asir and Mathematica general computer algebra system
for $\ensuremath{n=3}$ to $n=15$.

\parbox[t]{1\columnwidth}{%
\begin{center}
Table 2. Number of Independent Derivation Relations for MZVs
\par\end{center}
\begin{center}
\begin{tabular}{|c|c|c|c|c|c|c|c|c|c|c|c|c|c|c|c|}
\hline 
{\small{}$\vphantom{\dfrac{A}{B}}n$} & {\small{}3} & {\small{}4} & {\small{}5} & {\small{}6} & {\small{}7} & {\small{}8} & {\small{}9} & {\small{}10} & {\small{}11} & {\small{}12} & {\small{}13} & {\small{}14} & {\small{}15} & \textcolor{gray}{\small{}16{*}} & \textcolor{gray}{\small{}17{*}}\tabularnewline
\hline 
{\small{}$\vphantom{\dfrac{A}{B}}\mathscr{L}(n)$} & {\small{}1} & {\small{}2} & {\small{}5} & {\small{}10} & {\small{}22} & {\small{}44} & {\small{}90} & {\small{}181} & {\small{}363} & {\small{}727} & {\small{}1456} & {\small{}2912} & {\small{}5825} & \textcolor{gray}{\small{}11652{*}} & \textcolor{gray}{\small{}23304{*}}\tabularnewline
\hline 
\end{tabular}
\par\end{center}%
}

{*}{\small{} Calculated with a JavaScript application of my own development.}\footnote{\begin{onehalfspace}
On page https://tamas-tech.github.io 
\end{onehalfspace}
} 

\begin{minipage}[t]{1\columnwidth}%
\begin{shaded}%
\begin{conjecture}
For all weight $n\geq3$ the number of linearly independent relations
of Ihara-Kaneko-Zagier derivation relation can be calculated by the
following Euler-type pentagonal formula:
\[
\mathscr{L}(n)=2^{n-2-P_{1}}+2^{n-2-P_{2}}-2^{n-2-P_{3}}-2^{n-2-P_{4}}\pm\cdots-\omega(n-1)-1=-\sum_{k=1}^{n-2}\omega(k)\,2^{n-2-k}-\omega(n-1)-1.
\]
 The generating function of $\mathscr{L}(n)$ is given by: 
\[
{\displaystyle F(x)=\frac{x^{2}}{x-1}+\frac{x(x-1)}{2x-1}\prod_{k\geq1}(1-x^{k})},
\]
and $\mathscr{L}(n)$ satisfies the following recurrence relation
for all $n\geq3$:
\[
\mathscr{L}(n+1)=2\cdot\mathscr{L}(n)+1-\omega(n)+\omega(n-1).
\]
\end{conjecture}

\end{shaded}%
\end{minipage}

The calculations below verify that Conjecture 1. holds for $n=3$
through $n=17$:
\begin{flushleft}
$\mathscr{L}(3)=2^{1-{\color{blue}1}}-{\color{red}\omega(2)}-1=2^{0}+{\color{red}1}-1=1+1-1=\mathbf{1}$
\par\end{flushleft}

\begin{flushleft}
$\mathscr{L}(4)=2^{2-{\color{blue}1}}+2^{2-{\color{blue}2}}-{\color{red}\omega(3)}-1=2^{1}+2^{0}-{\color{red}0}-1=2+1-1=\mathbf{2}$
\par\end{flushleft}

\begin{flushleft}
$\mathscr{L}(5)=2^{3-{\color{blue}1}}+2^{3-{\color{blue}2}}-{\color{red}\omega(4)}-1=2^{2}+2^{1}-{\color{red}0}-1=4+2-1=\mathbf{5}$
\par\end{flushleft}

\begin{flushleft}
$\mathscr{L}(6)=2^{4-{\color{blue}1}}+2^{4-{\color{blue}2}}-{\color{red}\omega(5)}-1=2^{3}+2^{2}-{\color{red}1}-1=8+4-1-1=\mathbf{10}$
\par\end{flushleft}

\begin{flushleft}
$\mathscr{L}(7)=2^{5-{\color{blue}1}}+2^{5-{\color{blue}2}}-2^{5-{\color{blue}5}}-{\color{red}\omega(6)}-1=2^{4}+2^{3}-2^{0}-{\color{red}0}-1=16+8-1-1=\mathbf{22}$
\par\end{flushleft}

\begin{flushleft}
$\mathscr{L}(8)=2^{6-{\color{blue}1}}+2^{6-{\color{blue}2}}-2^{6-{\color{blue}5}}-{\color{red}\omega(7)}-1=2^{5}+2^{4}-2^{1}-{\color{red}1}-1=32+16-2-1-1=\mathbf{44}$
\par\end{flushleft}

\begin{flushleft}
$\mathscr{L}(9)=2^{7-{\color{blue}1}}+2^{7-{\color{blue}2}}-2^{7-{\color{blue}5}}-2^{7-{\color{blue}7}}-{\color{red}\omega(8)}-1=2^{6}+2^{5}-2^{2}-2^{0}-{\color{red}0}-1=64+32-4-1-1=\mathbf{90}$
\par\end{flushleft}

\begin{flushleft}
$\mathscr{L}(10)=2^{8-{\color{blue}1}}+2^{8-{\color{blue}2}}-2^{8-{\color{blue}5}}-2^{8-{\color{blue}7}}-{\color{red}\omega(9)}-1=2^{7}+2^{6}-2^{3}-2^{1}-{\color{red}0}-1=128+64-8-2-1=\mathbf{181}$
\par\end{flushleft}

\begin{flushleft}
$\mathscr{L}(11)=2^{9-{\color{blue}1}}+2^{9-{\color{blue}2}}-2^{9-{\color{blue}5}}-2^{9-{\color{blue}7}}-{\color{red}\omega(10)}-1=2^{8}+2^{7}-2^{4}-2^{2}-{\color{red}0}-1=256+128-16-4-1=\mathbf{363}$
\par\end{flushleft}

\begin{flushleft}
$\mathscr{L}(12)=2^{10-{\color{blue}1}}+2^{10-{\color{blue}2}}-2^{10-{\color{blue}5}}-2^{10-{\color{blue}7}}-{\color{red}\omega(11)}-1=2^{9}+2^{8}-2^{5}-2^{3}-{\color{red}0}-1=512+256-32-8-1=\mathbf{727}$
\par\end{flushleft}

\begin{flushleft}
$\mathscr{L}(13)=2^{11-{\color{blue}1}}+2^{11-{\color{blue}2}}-2^{11-{\color{blue}5}}-2^{11-{\color{blue}7}}-{\color{red}\omega(12)}-1=2^{10}+2^{9}-2^{6}-2^{4}+{\color{red}1}-1=1024+512-64-16=\mathbf{1456}$
\par\end{flushleft}

\begin{flushleft}
$\mathscr{L}(14)=2^{12-{\color{blue}1}}+2^{12-{\color{blue}2}}-2^{12-{\color{blue}5}}-2^{12-{\color{blue}7}}+2^{12-{\color{blue}12}}-{\color{red}\omega(13)}-1=2^{11}+2^{10}-2^{7}-2^{5}+2^{0}-{\color{red}0}-1=2048+1024-128-32+1-1=\mathbf{2912}$
\par\end{flushleft}

\begin{flushleft}
$\mathscr{L}(15)=2^{13-{\color{blue}1}}+2^{13-{\color{blue}2}}-2^{13-{\color{blue}5}}-2^{13-{\color{blue}7}}+2^{13-{\color{blue}12}}-{\color{red}\omega(14)}-1=2^{12}+2^{11}-2^{8}-2^{6}+2^{1}-{\color{red}0}-1=4096+2048-256-64+2-1=\mathbf{5825}$
\par\end{flushleft}

\begin{flushleft}
$\mathscr{L}(16)=2^{14-{\color{blue}1}}+2^{14-{\color{blue}2}}-2^{14-{\color{blue}5}}-2^{14-{\color{blue}7}}+2^{14-{\color{blue}12}}-{\color{red}\omega(15)}-1=2^{13}+2^{12}-2^{9}-2^{7}+2^{2}+{\color{red}1}-1=8192+4096-512-128+4=\mathbf{11652}$
\par\end{flushleft}

\begin{flushleft}
$\mathscr{L}(17)=2^{15-{\color{blue}1}}+2^{15-{\color{blue}2}}-2^{15-{\color{blue}5}}-2^{15-{\color{blue}7}}+2^{15-{\color{blue}12}}+2^{15-{\color{blue}15}}-{\color{red}\omega(16)}-1=2^{14}+2^{13}-2^{10}-2^{8}+2^{3}+2^{0}-{\color{red}0}-1=$
\par\end{flushleft}

$\hphantom{\mathscr{L}(17)}=16384+8192-1024-256+8+1-1=\mathbf{23304}$

~

The first equivalent version reveals a simple connection between $\mathscr{L}(n)$
and $\text{pic}(n)=$\href{https://oeis.org/A178841}{A178841(n)}
the number of pure inverting compositions of $n$. 
\begin{conjecture}
For all weight $n\geq3$ the number of linearly independent relations
of Ihara-Kaneko-Zagier derivation relation can be calculated by the
following formula:
\[
\mathscr{L}(n)=2^{n-2}-1-\text{pic}(n-1)
\]
\end{conjecture}

The calculations summarized in the table below verify that Conjecture
2. holds for $n=3$ through $n=17$:

\parbox[t]{1\columnwidth}{%
\begin{center}
Table 3. $\mathscr{L}(n)=2^{n-2}-1-\text{pic}(n-1)$
\par\end{center}
\begin{center}
{\small{}}%
\begin{tabular}{|c||c|c|c|c|c|c|c|c|c|c|c|c|c|c|c|}
\hline 
{\small{}$\vphantom{\dfrac{A}{B}}n$} & {\small{}3} & {\small{}4} & {\small{}5} & {\small{}6} & {\small{}7} & {\small{}8} & {\small{}9} & {\small{}10} & {\small{}11} & {\small{}12} & {\small{}13} & {\small{}14} & {\small{}15} & {\small{}16} & {\small{}17}\tabularnewline
\hline 
{\small{}$\vphantom{\dfrac{A}{B}}2^{n-2}-1$} & {\small{}1} & {\small{}3} & {\small{}7} & {\small{}15} & {\small{}31} & {\small{}63} & {\small{}127} & {\small{}255} & {\small{}511} & {\small{}1023} & {\small{}2047} & {\small{}4095} & {\small{}8191} & {\small{}16383} & {\small{}32767}\tabularnewline
\hline 
{\small{}$\vphantom{\dfrac{A}{B}}\text{pic}(n-1)$} & {\small{}0} & {\small{}1} & {\small{}2} & {\small{}5} & {\small{}9} & {\small{}19} & {\small{}37} & {\small{}74} & {\small{}148} & {\small{}296} & {\small{}591} & {\small{}1183} & {\small{}2366} & {\small{}4731} & {\small{}9463}\tabularnewline
\hline 
{\small{}$\vphantom{\dfrac{A}{B}}2^{n-2}-1-\text{pic}(n-1)$} & 1 & 2 & 5 & 10 & 22 & 44 & 90 & 181 & 363 & 727 & 1456 & 2912 & 5825 & 11652 & 23304\tabularnewline
\hline 
\end{tabular}{\small\par}
\par\end{center}%
}

Let $c(n)=$\href{https://oeis.org/A152537}{A152537(n)}. The second
equivalent reformulation highlights a deeper algebraic relationship
between $c(n)$ and $\text{pic}(n)$.
\begin{conjecture}
For all weight $n\geq3$ the number of linearly independent relations
of Ihara-Kaneko-Zagier derivation relation can be calculated by the
following formula:
\[
\mathscr{L}(n)=2^{n-2}-1-\omega(n-1)-c(n-2)
\]
\end{conjecture}

The calculations summarized in the table below verify that Conjecture
3. holds for $n=3$ through $n=17$:

\parbox[t]{1\columnwidth}{%
\begin{center}
Table 4. $\mathscr{L}(n)=2^{n-2}-1-\omega(n-1)-c(n-2)$
\par\end{center}
\begin{center}
{\small{}}%
\begin{tabular}{|c||c|c|c|c|c|c|c|c|c|c|c|c|c|c|c|}
\hline 
{\small{}$\vphantom{\dfrac{A}{B}}n$} & {\small{}3} & {\small{}4} & {\small{}5} & {\small{}6} & {\small{}7} & {\small{}8} & {\small{}9} & {\small{}10} & {\small{}11} & {\small{}12} & {\small{}13} & {\small{}14} & {\small{}15} & {\small{}16} & {\small{}17}\tabularnewline
\hline 
{\small{}$\vphantom{\dfrac{A}{B}}2^{n-2}-1$} & {\small{}1} & {\small{}3} & {\small{}7} & {\small{}15} & {\small{}31} & {\small{}63} & {\small{}127} & {\small{}255} & {\small{}511} & {\small{}1023} & {\small{}2047} & {\small{}4095} & {\small{}8191} & {\small{}16383} & {\small{}32767}\tabularnewline
\hline 
{\small{}$\vphantom{\dfrac{A}{B}}\omega(n-1)$} & {\small{}-1} & {\small{}0} & {\small{}0} & {\small{}1} & {\small{}0} & {\small{}1} & {\small{}0} & {\small{}0} & {\small{}0} & {\small{}0} & {\small{}-1} & {\small{}0} & {\small{}0} & {\small{}-1} & {\small{}0}\tabularnewline
\hline 
{\small{}$\vphantom{\dfrac{A}{B}}c(n-2)$} & {\small{}1} & {\small{}1} & {\small{}2} & {\small{}4} & {\small{}9} & {\small{}18} & {\small{}37} & {\small{}74} & {\small{}148} & {\small{}296} & {\small{}592} & {\small{}1183} & {\small{}2366} & {\small{}4732} & {\small{}9463}\tabularnewline
\hline 
{\small{}$\vphantom{\dfrac{A}{B}}2^{n-2}-1-\omega(n-1)-c(n-2)$} & 1 & 2 & 5 & 10 & 22 & 44 & 90 & 181 & 363 & 727 & 1456 & 2912 & 5825 & 11652 & 23304\tabularnewline
\hline 
\end{tabular}{\small\par}
\par\end{center}%
}

\section{Empirical and Theoretical Results Leading to Conjecture 1}

Before proceeding, we establish some terminology and notation that
will be used throughout the remainder of this work.

Le $\mathfrak{H}_{n}:=\{w\in\mathfrak{H},\left|w\right|=n\}$ denote
the set of words of weight $n$. We recall that the set of admissible
words is given by $\mathcal{A}=x\mathfrak{H}y$, and we denote the
subset of admissible words of weight $n$ by $\mathcal{A}_{n}:=x\mathfrak{H}_{n-2}y$.
For any admissible word $w\in\mathcal{A}$, we define its binary value,
denoted by $\mathrm{bin}(w)$ , as the unique even natural number
whose binary representation yields the word $w$ upon replacing each
digit $1$ with $x$ and each digit $0$ with $y$. Equivalently,
if a word is written as $\ensuremath{w=u_{1}u_{2}\dots u_{k}}$, where
$\ensuremath{u_{i}\in\{x,y\}}$ then 
\[
\ensuremath{\mathrm{bin}(w)=\sum_{i=1}^{k}c_{i}2^{k-i}},\text{ where }\ensuremath{c_{i}=1}\text{ if }\ensuremath{u_{i}=x}\text{, and }\ensuremath{c_{i}=0}\text{ if }\ensuremath{u_{i}=y.}
\]

For example, consider the admissible word $w=xxyy$. Replacing each
$x$ with $1$ and $y$ with $0$ gives the binary string $1100_{2}$.
Converting this to base 10, we obtain: $\mathrm{bin}(xxyy)=1\cdot2^{3}+1\cdot2^{2}+0\cdot2^{1}+0\cdot2^{0}=8+4=12.$

We will use a graded lexicographic total order on the words of $\mathfrak{H}$,
defined as follows. If $u$ and $v$ are words with $\left|u\right|<\left|v\right|$
, set $u<v$ . If $\left|u\right|=\left|v\right|$ and $u=wyw_{2}$
and $v=wxw_{3}$ , for any words $w$,$w_{1}$, and $w_{2}$, set
$u<v$ . It is easy to see that the order on $\mathcal{A}$ induced
by the mapping $\mathrm{bin(\,\cdot\,)}$ coincides with the graded
lexicographic order restricted to $\mathcal{A}$. The words of degree
at most $3$ are thus ordered as follows:
\[
1<y<x<yy<yx<xy<xx<yyy<yyx<yxy<yxx<xyy<xyx<xxy<xxx
\]

For any subset $V\subseteq\mathfrak{H}$ and any integer $m\ge1$,
let $\partial_{m}(V)$ denote the set $\{\partial_{m}(w):w\in V\}$. 

Conjecture 1 is based on the following analysis of the rank growth,
which chronologically preceded the formulation of Conjecture~2.

For a fixed $n\ge3$, we will construct a full row-rank matrix $\mathbf{M}_{n}$
from the derivatives $\partial_{n-\vert w\vert}(w)$ of the elements
$w$ in the subset
\[
\mathcal{B}_{n}=\mathcal{A}_{n-1}\cup x\mathcal{A}_{n-3}\cup x\mathcal{A}_{n-2}\cup\dots\cup x\mathcal{A}_{2}\cup\mathcal{A}_{2}\subset h
\]
$\left(\left|\mathcal{B}_{n}\right|=2^{n-3}+(2^{n-5}+2^{n-6}+\cdots+2^{0})+1=3\cdot2^{n-4}\right)$
according to the following procedure: 

The columns of $\mathbf{M}_{n}$ are indexed by the elements of the
set of admissible words of $\mathcal{A}_{n}$, ordered ascendingly
with respect to the graded lexicographic order. 
\begin{enumerate}
\item Let $w=x^{n-2}y$ be the maximal element of the set $\mathcal{A}_{n-1}$.
The first row of the matrix $\mathbf{M}_{n}$ is defined as the coordinate
vector of the derivative $\partial_{1}(w)\in\partial_{1}(\mathcal{A}_{n-1})\subset\text{Span}(\mathcal{A}_{n})$
relative to the ordered basis $(\mathcal{A}_{n},<)$.
\item We then proceed iteratively through the remaining elements of $A_{n-1}$
in descending order. For each element, we compute the coordinate vector
of its derivative and insert it as a new row in $\mathbf{M}_{n}$.
\emph{Crucially, throughout this process, we ensure that if a candidate
row can be expressed as a linear combination of the preceding rows,
it is discarded}.
\item This procedure is then repeated analogously for the elements of the
remaining subsets in $\mathcal{B}_{n}$ (namely $x\mathcal{A}_{n-3},x\mathcal{A}_{n-2},\dots,x\mathcal{A}_{2}$,
and $\mathcal{A}_{2}$ in the given order), examining their corresponding
derivatives in descending order within each subset. 
\end{enumerate}
Constructing the matrix via this elimination process, it is evident
that the resulting matrix $\mathbf{M}_{n}$ will have full row-rank.
At the end of the procedure, we count how many derivatives originating
from the subsets $\mathcal{A}_{n-1},x\mathcal{A}_{n-3},x\mathcal{A}_{n-2},\dots,x\mathcal{A}_{2}$,
and $\mathcal{A}_{2}$ remain as rows in the matrix. These quantities
will be denoted by $r_{0,n-2},r_{1,n-2},r_{2,n-2},\dots,r_{n-3,n-2}$
respectively, and they will constitute the entries of the $(n-2)$-th
column of an upper triangular matrix $\mathbf{R}$, whose rows and
columns are indexed starting from $0$.

Since our JavaScript application, which includes the implementation
of the above procedure, generates a clear graphical layout for smaller
values of $n$, the visual output representing the matrix $\mathbf{M}_{6}$
for $n=6$ is included below as an illustrative example.
\begin{center}
\includegraphics[width=0.96\textwidth]{Rank6}
\par\end{center}

For $n=6$, the generating set is given by $\mathcal{B}_{6}=\mathcal{A}_{5}\cup x\mathcal{A}_{3}\cup x\mathcal{A}_{2}\cup\mathcal{A}_{2}$.
As clearly observed from the visual output, all $2^{3}=8$ elements
of $\partial_{1}\left(\mathcal{A}_{5}\right)$ and all $2^{1}=2$
elements of $\partial_{2}\left(x\mathcal{A}_{3}\right)$ remain in
the matrix. Conversely, the elements of $\partial_{3}\left(x\mathcal{A}_{2}\right)=\{\partial_{3}(xxy)\}$
and $\partial_{4}\left(\mathcal{A}_{2}\right)=\{\partial_{4}(xy)\}$
can already be expressed as linear combinations of the preceding rows
and are thus discarded. By definition, this yields $\ensuremath{r_{0,4}=8},\,\ensuremath{r_{1,4}=2}$,
and $r_{2,4}=r_{3,4}=0$.

To provide an additional illustration, let us consider the output
generated for the case $n=7$.
\begin{center}
\includegraphics[width=0.96\textwidth]{Rank7}
\par\end{center}

From this output, it can be observed that all $16$, $4$, and $2$
elements of $\partial_{1}\left(\mathcal{A}_{6}\right)$, $\partial_{2}\left(x\mathcal{A}_{4}\right)$,
and $\partial_{3}\left(x\mathcal{A}_{3}\right)$, respectively, remain
in the matrix, whereas all elements of $\partial_{4}\left(x\mathcal{A}_{2}\right)$
and $\partial_{5}\left(\mathcal{A}_{2}\right)$ are eliminated. By
definition, this yields $r_{0,5}=16$, $r_{1,5}=4$, $r_{2,5}=2$,
and $r_{3,5}=r_{4,5}=0$.

Next, we present the empirical results obtained for $r_{i,j}$ from
$n=2$ to $n=17$, collected into the table below:

\[
\begin{array}{cc|ccccccccccccccccc}
{\color{lightgray}} & {\color{lightgray}n=} & {\color{lightgray}2} & {\color{lightgray}3} & {\color{lightgray}4} & {\color{lightgray}5} & {\color{lightgray}6} & {\color{lightgray}7} & {\color{lightgray}8} & {\color{lightgray}9} & {\color{lightgray}10} & {\color{lightgray}11} & {\color{lightgray}12} & {\color{lightgray}13} & {\color{lightgray}14} & {\color{lightgray}15} & {\color{lightgray}16} & {\color{lightgray}17} & {\color{lightgray}18}\\
{\color{lightgray}} & i\setminus j & 0 & 1 & 2 & 3 & 4 & 5 & 6 & 7 & 8 & 9 & 10 & 11 & 12 & 13 & 14 & 15 & 16\\
\hline {\color{lightgray}\partial_{1}} & 0 & 1 & 1 & 2 & 4 & 8 & 16 & 32 & 64 & 128 & 256 & 512 & 1024 & 2048 & 4096 & 8192 & 16384 & \cdots\\
{\color{lightgray}\partial_{2}} & 1 &  &  & 0 & 1 & 2 & 4 & 8 & 16 & 32 & 64 & 128 & 256 & 512 & 1024 & 2048 & 4096 & \cdots\\
{\color{lightgray}\partial_{3}} & 2 &  &  &  & 0 & 0 & 2 & 3 & 6 & 12 & 24 & 48 & 96 & 192 & 384 & 768 & 1536 & \cdots\\
{\color{lightgray}\partial_{4}} & 3 &  &  &  &  & 0 & 0 & 1 & 3 & 6 & 10 & 21 & 42 & 84 & 168 & 336 & 672 & \cdots\\
{\color{lightgray}\partial_{5}} & 4 &  &  &  &  &  & 0 & 0 & 1 & 2 & 6 & 10 & 20 & 39 & 78 & 158 & 315 & \cdots\\
{\color{lightgray}\partial_{6}} & 5 &  &  &  &  &  &  & 0 & 0 & 1 & 2 & 5 & 10 & 20 & 38 & 76 & 152 & \cdots\\
{\color{lightgray}\partial_{7}} & 6 &  &  &  &  &  &  &  & 0 & 0 & 1 & 2 & 5 & 9 & 20 & 38 & 75 & \cdots\\
{\color{lightgray}\partial_{8}} & 7 &  &  &  &  &  &  &  &  & 0 & 0 & 1 & 2 & 5 & 9 & 19 & 38 & \cdots\\
{\color{lightgray}\partial_{9}} & 8 &  &  &  &  &  &  &  &  &  & 0 & 0 & 1 & 2 & 5 & 9 & 19 & \cdots\\
{\color{lightgray}\partial_{10}} & 9 &  &  &  &  &  &  &  &  &  &  & 0 & 0 & 1 & 2 & 5 & 9 & \cdots\\
{\color{lightgray}\partial_{11}} & 10 &  &  &  &  &  &  &  &  &  &  &  & 0 & 0 & 1 & 2 & 5 & \cdots\\
{\color{lightgray}\partial_{12}} & 11 &  &  &  &  &  &  &  &  &  &  &  &  & 0 & 0 & 1 & 2 & \cdots\\
{\color{lightgray}\partial_{13}} & 12 &  &  &  &  &  &  &  &  &  &  &  &  &  & 0 & 0 & 1 & \cdots\\
{\color{lightgray}\partial_{14}} & 13 &  &  &  &  &  &  &  &  &  &  &  &  &  &  & 0 & 0 & \cdots\\
{\color{lightgray}\partial_{15}} & 14 &  &  &  &  &  &  &  &  &  &  &  &  &  &  &  & 0 & \cdots\\
{\color{lightgray}\partial_{16}} & 15 &  &  &  &  &  &  &  &  &  &  &  &  &  &  &  &  & \cdots\cdots
\end{array}
\]

The most significant observation regarding the resulting table is
that, putting aside the first few entries of each row, the values
follow a clearly discernible recurrence relation. Specifically, the
$j$-th entry of the $i$-th row can be obtained as the difference
between the $(j-1)$-th entry of the $(i-1)$-th row and the $(j-i-1)$-th
entry of the $(i-1)$-th row. 

Formally, for sufficiently large indices, the entries satisfy the
following recurrence: $r_{i,j}=r_{i-1,j-1}-r_{i-1,j-i-1}$. To illustrate
this pattern, consider the following examples taken directly from
the empirical data:
\begin{itemize}
\item For $i=4$ and $j=12$ (corresponding to $\partial_{5}$ at $n=14$),
we have $39=42-3$.
\item For $i=3$ and $j=13$ (corresponding to $\partial_{4}$ at $n=15$),
we have $168=192-24$.
\end{itemize}
By applying this recurrence relation backward, we can determine the
exact initial values required at the beginning of each row for the
relation to hold globally across the entire table. Although such a
backward extension could potentially lead to a contradiction, remarkably,
no such inconsistency arises. Specifically, if we set all entries
on the main diagonal to $1$ and all entries strictly below the main
diagonal to $0$, the recurrence relation is satisfied consistently
throughout the entire table. The resulting extended table is presented
below:
\[
\begin{array}{c|ccccccccccccccccc}
i\setminus j & 0 & 1 & 2 & 3 & 4 & 5 & 6 & 7 & 8 & 9 & 10 & 11 & 12 & 13 & 14 & 15 & 16\\
\hline 0 & \mathbf{1} & \mathbf{1} & \mathbf{2} & \mathbf{4} & \mathbf{8} & \mathbf{16} & \mathbf{32} & \mathbf{64} & \mathbf{128} & \mathbf{256} & \mathbf{512} & \mathbf{1024} & \mathbf{2048} & \mathbf{4096} & \mathbf{8192} & \mathbf{16384} & \cdots\\
1 & 0 & 1 & 0 & 1 & 2 & 4 & 8 & 16 & 32 & 64 & 128 & 256 & 512 & 1024 & 2048 & 4096 & \cdots\\
2 & 0 & 0 & 1 & 0 & 0 & 2 & 3 & 6 & 12 & 24 & 48 & 96 & 192 & 384 & 768 & 1536 & \cdots\\
3 & 0 & 0 & 0 & 1 & 0 & 0 & 1 & 3 & 6 & 10 & 21 & 42 & 84 & 168 & 336 & 672 & \cdots\\
4 & 0 & 0 & 0 & 0 & 1 & 0 & 0 & 1 & 2 & 6 & 10 & 20 & 39 & 78 & 158 & 315 & \cdots\\
5 & 0 & 0 & 0 & 0 & 0 & 1 & 0 & 0 & 1 & 2 & 5 & 10 & 20 & 38 & 76 & 152 & \cdots\\
6 & 0 & 0 & 0 & 0 & 0 & 0 & 1 & 0 & 0 & 1 & 2 & 5 & 9 & 20 & 38 & 75 & \cdots\\
7 & 0 & 0 & 0 & 0 & 0 & 0 & 0 & 1 & 0 & 0 & 1 & 2 & 5 & 9 & 19 & 38 & \cdots\\
8 & 0 & 0 & 0 & 0 & 0 & 0 & 0 & 0 & 1 & 0 & 0 & 1 & 2 & 5 & 9 & 19 & \cdots\\
9 & 0 & 0 & 0 & 0 & 0 & 0 & 0 & 0 & 0 & 1 & 0 & 0 & 1 & 2 & 5 & 9 & \cdots\\
10 & 0 & 0 & 0 & 0 & 0 & 0 & 0 & 0 & 0 & 0 & 1 & 0 & 0 & 1 & 2 & 5 & \cdots\\
11 & 0 & 0 & 0 & 0 & 0 & 0 & 0 & 0 & 0 & 0 & 0 & 1 & 0 & 0 & 1 & 2 & \cdots\\
12 & 0 & 0 & 0 & 0 & 0 & 0 & 0 & 0 & 0 & 0 & 0 & 0 & 1 & 0 & 0 & 1 & \cdots\\
13 & 0 & 0 & 0 & 0 & 0 & 0 & 0 & 0 & 0 & 0 & 0 & 0 & 0 & 1 & 0 & 0 & \cdots\\
14 & 0 & 0 & 0 & 0 & 0 & 0 & 0 & 0 & 0 & 0 & 0 & 0 & 0 & 0 & 1 & 0 & \cdots\\
15 & 0 & 0 & 0 & 0 & 0 & 0 & 0 & 0 & 0 & 0 & 0 & 0 & 0 & 0 & 0 & 1 & \cdots\\
\vdots & \vdots & \vdots & \vdots & \vdots & \vdots & \vdots & \vdots & \vdots & \vdots & \vdots & \vdots & \vdots & \vdots & \vdots & \vdots & \vdots & \ddots\\
{\displaystyle \sum_{j=1}^{i}a_{i,j}} & 1 & \mathbf{2} & \mathbf{3} & \mathbf{6} & \boldsymbol{11} & \mathbf{23} & \mathbf{45} & \boldsymbol{91} & \mathbf{182} & \mathbf{364} & \mathbf{728} & \mathbf{1457} & \mathbf{2913} & \mathbf{5826} & \mathbf{11653} & \mathbf{23305} & \cdots
\end{array}
\]

It naturally follows from the construction of the table that subtracting
$1$ from the sum of the entries in its $j$-th column yields exactly
the values of $\mathscr{L}(j+2)$ for all $j\geq1$. 

A formalization of the recurrence relation observed in the table is
given by:

\[
r_{i,j}=\begin{cases}
0 & \text{if }i<0\text{ or }j<0;\\
1 & \text{if }i=j\text{ for }i=0,1,2,3,\dots;\\
2^{\max(0,j-1)} & \text{if }i=1\text{ for }j=0,1,2,3,\dots;\\
a_{i-1,j-1}-a_{i-1,j-i-1} & \text{otherwise.}
\end{cases}
\]

The recurrence relation was investigated within a more general theoretical
framework as a transformation. Instead of the initial values $2^{\max(0,j-1)}$
entered into the first row, let us consider an arbitrary arithmetic
sequence $\boldsymbol{a}=(a(1),a(2),\dots)$, and let the column sums
define the values of the transformed sequence $T(\boldsymbol{a})=(T(a(1)),T(a(2)),\dots)$.
During our thorough theoretical investigation, it soon became apparent
that this transformation is remarkably versatile and well-suited for
providing concise and clear explanations of the connections between
various combinatorial results with a special emphasis on pentagonal
formulas. Furthermore, since it is based on shifted differences and
a summation, the transformation is inherently linear. What is essential
for our purposes is that the transformation maps the sequence $\boldsymbol{a}=(1,0,0,0,\dots)$
to the sequence $-\boldsymbol{\omega}=(1,1,0,0,-1,0,-1,0,0,0,0,1,0,\dots)$
(now, starting from $1$ rather than $0$ as done above). Consequently,
by linearity, the transform of an arbitrary sequence $\boldsymbol{a}=(a_{1},a_{2},a_{3},\dots)$
yields a pentagonal formula of the form $-\sum_{k=1}^{\infty}\omega(k)a_{n-k}$,
which can be written more concisely as $T(\boldsymbol{a})=-(a_{1},a_{2},\dots)*(\omega(1),\omega(2),\dots)$,
where $*$ denotes the convolution of sequences. This, along with
a few straightforward algebraic manipulations, directly recovers the
pentagonal formula presented in Conjecture 1. Since the specific areas
discussing this type of transformation lie beyond the current scope
of our knowledge, we shall adopt the term \emph{pentagonal transformation}
for the remainder of this work. During the study of the transformation,
it also became apparent that when transforming various sequences $\boldsymbol{a}=(a(1),a(2),\dots)$,
it is not only the transformed sequence $T(\mathbf{a})$ itself that
is of interest; the individual rows generated in the table during
the underlying recurrence relation, as well as their corresponding
generating functions, are also highly significant. In the table above
leading to the values of $\mathscr{L}(j)$, for instance, it is clearly
observable that the first non-zero entries of successive rows coincide
with the sequence $(\text{pic}(0),\text{pic}(1),\text{pic}(2),\dots)=(1,0,0,1,2,5,9,19,37,\dots)$
for increasingly longer stretches. From this, together with certain
relations to be discussed in Section 5, it will become completely
evident that each row of the table contains parametric generalizations
$(\text{pic}_{n}(0),\text{pic}_{n}(1),\text{pic}_{n}(2),\dots)$ of
the sequence $(\text{pic}(0),\text{pic}(1),\text{pic}(2),\dots)$
corresponding to the parameter $n$ of that row. 

As for the proof of Conjecture 1, it suffices to demonstrate two conditions:
\begin{enumerate}
\item for any $n\geq3$, the vectors in the set $\partial_{1}(\mathcal{A}_{n-1})$
are linearly independent; 
\item the rank increments $r_{i,j}$ follow the aforementioned recurrence
relation. 
\end{enumerate}
The first statement can be straightforwardly established without any
difficulty, from which the corresponding power of two naturally follows.
However, we have been unable to provide a proof for the second statement.

It is worth emphasizing that the credibility of Conjecture 1 is significantly
reinforced by the fact that the consistency does not merely rest on
the first fourteen terms of a single sequence. Rather, it is deeply
embedded in the entire underlying table of rank increments, matching
perfectly across the whole two-dimensional recursive structure. 

\section{Empirical and Theoretical Results Leading to Conjecture 2}

The empirical results leading to Conjecture 2 can be obtained via
the exact same elimination procedure described for Conjecture~1;
however, rather than traversing the words in descending order, the
matrix is iteratively expanded by examining the coordinate vectors
of their derivatives in ascending order. This modification offers
a distinct computational advantage: as can be readily verified, the
coordinate vectors of the derivatives in the set 
\[
\mathcal{B}_{n}^{*}=\mathcal{B}_{n}\setminus\mathcal{A}_{n-1}=\partial_{n-2}(\mathcal{A}_{2})\cup\partial_{n-3}(x\mathcal{A}_{2})\cup\dots\cup\partial_{2}(x\mathcal{A}_{n-3})
\]
relative to the ordered basis $(\mathcal{A}_{n},<)$ are guaranteed
to form a linearly independent system of vectors. Consequently, for
any fixed $n\ge3$, the elimination procedure can be initiated immediately
with the full row-rank matrix $\mathbf{N}_{n}$ whose rows are defined
as the coordinate vectors of the derivatives contained in the set
$\mathcal{B}_{n}^{*}$ without requiring any prior rank computation.
This initial matrix $\mathbf{N}_{n}$ consists of exactly $1+(2^{0}+2^{1}+\cdots+2^{n-6}+2^{n-5})=2^{n-4}$
rows. Subsequently, we iterate through the remaining elements of the
set $\mathcal{A}_{n-1}$ in ascending order, starting from its minimal
element, $xy^{n-2}$. For each successive word $w$, we append the
coordinate vector of its derivative $\partial_{1}(w)$ as a candidate
row to $\mathbf{N}_{n}$; if this candidate row increases the rank
of the matrix, it is retained, and otherwise, it is discarded. Upon
completion of this procedure, the resulting matrix $\mathbf{N}_{n}$
is inherently full row-rank, just like the matrix $\mathbf{M}_{n}$
discussed in Conjecture 1. Furthermore, the ranks of these two matrices
are identical. 

Since our custom JavaScript application—which implements the aforementioned
procedure too—generates a clear graphical layout for small values
of $n$, the visual outputs representing the matrices $\mathbf{N}_{6}$
and $\mathbf{N}_{7}$ are included below as illustrative examples. 
\begin{center}
\includegraphics[width=0.96\textwidth]{Rank6Le}
\par\end{center}

\begin{center}
\includegraphics[width=0.96\textwidth]{Rank7Le}
\par\end{center}

Although it is theoretically possible that a word $w_{1}\in\mathcal{A}_{n-1}$
is skipped while the derivative of a larger word $w_{2}>w_{1}$ is
retained in the matrix $\mathbf{N}_{n}$, empirical evidence suggests
that this never occurs for $n>6$. Specifically, for any weight $n\neq6$,
there exists a unique critical threshold word $w_{\text{crit}}\in\mathcal{A}_{n-1}$
such that its derivative is discarded, and consequently, the derivatives
of all subsequent words $w>w_{\text{crit}}$ are also eliminated.
In short, with the sole exception of $n=6$, elements are systematically
discarded exclusively from the tail of the ordered set $\mathcal{A}_{n-1}$.
(The anomalous behavior observed in the case $n=6$ is explicitly
accounted for by the \textquotedbl congestion\textquotedbl{} of the
pentagonal numbers 3 and 5.)

In light of these observations, \emph{Conjecture 2 can be stated in
a sharper form as follows}: for any weight $n\neq6$, during the construction
of the matrix $\mathbf{N}_{n}$, the first and minimal word $w_{0}\in\mathcal{A}_{n-1}$
whose derivative fails to increase the rank of the matrix is precisely
the $\left(3\cdot2^{n-4}-\text{pic}(n-1)\right)$-th element in the
ordered set. Consequently, the derivatives of all words larger than
$w_{0}$ likewise fail to increment the rank of the matrix. 

Since $|\mathcal{B}_{n}^{*}|=2^{n-4}$, and exactly $3\cdot2^{n-4}-\text{pic}(n-1)-1$
elements from $\mathcal{A}_{n-1}$ are retained in the matrix $\mathbf{N}_{n}$,
it follows that: 
\[
\mathscr{L}(n)=\text{rank}(\mathbf{N}_{n})=2^{n-4}+3\cdot2^{n-4}-\text{pic}(n-1)-1=2^{n-2}-\text{pic}(n-1)-1,
\]
which precisely recovers the formulation of Conjecture 2 presented
in Section 2. Although the equivalence between Conjecture 1 and Conjecture
2 can be formally established with ease in light of the identities
to be discussed in Section 5 (where we shall demonstrate that the
rank increments are fundamentally parametric generalizations of pure
inverting compositions), we consider it essential to present these
two approaches separately, as the structural nuances arising between
them are not fully captured by the explicit formulas alone. 

It is worth noting that this particular formulation and matrix construction
are highly optimized for computationally efficient implementations,
which could potentially enable the verification of the conjecture
for weights well beyond $n=17$. 

It is highly instructive to reflect upon a comment by N. J. A. Sloane
in the OEIS entry for $\mathbf{A178841}$, where he noted that the
striking similarity between the terms of $\mathbf{A178841}(n)$ and
$\mathbf{A152537}(n)$ might be \textquotedbl just a coincidence,
caused by the fact that their generating functions are almost identical.\textquotedbl{}
Our framework provides a rigorous, deeper theoretical explanation
that dispels this assumption of coincidence. The connection between
these two sequences is fundamentally rooted in the dual matrix constructions
presented herein. Specifically, building the derivation matrix in
descending order yields a rank-increment structure governed by $\mathbf{A152537}$
(Conjecture 1), whereas building it in ascending order yields a structure
governed by $\mathbf{A178841}$ (Conjecture~2). Since both procedures
ultimately resolve the exact same system of Ihara-Kaneko\nobreakdash-Zagier
derivation relations, the ranks of the resulting matrices $\mathbf{M}_{n}$
and $\mathbf{N}_{n}$ must identically coincide with $\mathscr{L}(n)$,
thereby intrinsically forcing the near-identity of their respective
generating functions. 

\section{Generalization of $\text{pic}(n)$}
\begin{center}
\textcolor{gray}{Preliminary draft.}
\par\end{center}

In this section, we demonstrate that the rank increments can be viewed
as a natural generalization of the classical sequence $\text{pic}(k)$
to the parametric function $\text{pic}_{n}(k)$.

Before presenting the formal identities, we define three fundamental
families of polynomials that govern these combinatorial structures.
Let $\mathrm{Fib}_{n}$ and $\mathrm{Fab}_{n}$ denote the multivariate
Fibonacci and Faber partition polynomials, respectively. Furthermore,
let $\mathrm{Z}_{n}$ denote the cycle index polynomial of the symmetric
group $S_{n}$. For a given set of variables $(x_{1},x_{2},\ldots,x_{n})$,
these polynomials are explicitly defined as follows: 

\begin{eqnarray*}
\mathrm{Fib}_{n}(x_{1},x_{2},\ldots x_{n}) & = & \sum_{1b_{1}+2b_{2}+\cdots+nb_{n}=n}(-1)^{b_{1}+\cdots+b_{n}}\left(b_{1}+\cdots+b_{n}\right)!\;\frac{x_{1}^{b_{1}}x_{2}^{b_{2}}\cdots x_{n}^{b_{n}}}{b_{1}!\,b_{2}!\cdots b_{n}!}\\
\mathrm{Fab}_{n}(x_{1},x_{2},\ldots x_{n}) & = & \sum_{1b_{1}+2b_{2}+\cdots+nb_{n}=n}(-1)^{b_{1}+\cdots+b_{n}}\,n\,\left(b_{1}+\cdots+b_{n}-1\right)!\;\frac{x_{1}^{b_{1}}x_{2}^{b_{2}}\cdots x_{n}^{b_{n}}}{b_{1}!\,b_{2}!\cdots b_{n}!}\\
\mathrm{Z}_{n}(x_{1},x_{2},\ldots x_{n}) & = & \sum_{1b_{1}+2b_{2}+\cdots+nb_{n}=n}\dfrac{1}{1^{b_{1}}2^{b_{2}}\cdots n^{b_{n}}}\;\frac{x_{1}^{b_{1}}x_{2}^{b_{2}}\cdots x_{n}^{b_{n}}}{b_{1}!\,b_{2}!\cdots b_{n}!}
\end{eqnarray*}

Motivated by the theoretical results established in \cite{claesson2023permutationsinversions},
we present the fundamental recurrence relation satisfied by $\text{pic}(k)$,
along with several formulas that characterize the intricate relations
among the pure inverting composition function $\text{pic}(k)$, the
partition function $p(k)$, and the divisor sum function $\sigma(k)$. 

\begin{eqnarray*}
\text{pic}(k) & = & \sum_{j=0}^{k}\omega(j)\,2^{k-1-j}+\omega(j)\\
k\,\text{pic}(k) & = & \sum_{j=0}^{k}\text{pic}(k-j)\,\left(2^{j}-\sigma(j)-1\right)\\
\text{pic}(k) & = & 2\,\text{pic}(k-1)+\omega(k)-\omega(k-1)\\
\text{pic}(k) & = & \text{Z}_{k}(2^{1}-\sigma(1)-1,2^{2}-\sigma(2)-1,\ldots,2^{k}-\sigma(k)-1)\\
\sigma(k)-2^{k}+1 & = & \text{Fab}_{k}(\text{pic}(1),\text{pic}(2),\ldots,\text{pic}(k))\\
\sum_{j=0}^{k}p(j)-2p(k) & = & -\text{Fib}_{k}(\text{pic}(1),\text{pic}(2),\ldots,\text{pic}(k))\\
\text{pic}(k) & = & \text{Fib}_{k}\left(2p(1)-\sum_{j=0}^{1}p(j),2p(2)-\sum_{j=0}^{2}p(j),\ldots,2p(k)-\sum_{j=0}^{k}p(j)\right)\\
p(k) & = & \sum_{j=0}^{k}2^{\max(0,k-1-j)}\,\text{Fib}_{j}(\text{pic}(1),\text{pic}(2),\ldots,\text{pic}(j))
\end{eqnarray*}

These relationships can be naturally generalized by restricting the
underlying arithmetic functions. Specifically, let $p_{{\scriptscriptstyle >}n}(k)$
denote the number of partitions of $k$ into parts strictly greater
than $n$, and let $\sigma_{n}(k)$ denote the restricted divisor
sum function, representing the sum of divisors of $k$ that are at
most $n$. By adopting the first formula below as the formal definition
of the generalized function $\text{pic}_{n}(k)$, all the preceding
identities can be seamlessly mapped to this generalized setting as
follows: 

\begin{eqnarray*}
\text{pic}_{n}(k) & := & \text{Z}_{k}(2^{1}-\sigma_{n}(1)-1,2^{2}-\sigma_{n}(2)-1,\ldots,2^{k}-\sigma_{n}(k)-1)\\
k\,\text{pic}_{n}(k) & = & \sum_{j=1}^{k}\text{pic}_{n}(k-j)\,\left(2^{j}-\sigma_{n}(j)-1\right)\\
\text{pic}_{n}(k) & = & \text{pic}_{n-1}(k)-\text{pic}_{n-1}(k-n)\\
\text{pic}_{n}\left(\binom{n+1}{2}+k\right) & = & 2^{k-1}\left[2\,\text{pic}_{n}\left(\binom{n+1}{2}\right)-(-1)^{n}\right]\\
\sigma_{n}(k)-2^{k}+1 & = & \text{Fab}_{k}(\text{pic}_{n}(1),\text{pic}_{n}(2),\ldots,\text{pic}_{n}(k))\\
\sum_{j=0}^{k}p_{{\scriptscriptstyle >}n}(j)-2p_{{\scriptscriptstyle >}n}(k) & = & -\text{Fib}_{k}(\text{pic}_{n}(1),\text{pic}_{n}(2),\ldots,\text{pic}_{n}(k))\\
\text{pic}_{n}(k) & = & \text{Fib}_{k}\left(2p_{{\scriptscriptstyle >}n}(1)-\sum_{j=0}^{1}p_{{\scriptscriptstyle >}n}(j),2p_{{\scriptscriptstyle >}n}(2)-\sum_{j=0}^{2}p_{{\scriptscriptstyle >}n}(j),\ldots,2p_{{\scriptscriptstyle >}n}(k)-\sum_{j=0}^{k}p_{{\scriptscriptstyle >}n}(j)\right)\\
p_{{\scriptscriptstyle >}n}(k) & = & \sum_{j=0}^{k}2^{\max(0,k-1-j)}\,\text{Fib}_{j}(\text{pic}_{n}(1),\text{pic}_{n}(2),\ldots,\text{pic}_{n}(j))
\end{eqnarray*}

Furthermore, the following connection holds between the newly introduced
generalized function $\text{pic}_{n}(k)$ and the classical sequence
$\text{pic}(k)$: 
\[
\sum_{k=0}^{n}\text{pic}_{n-k}(k)=2^{n}-\text{pic}(n+1)\,\,(n=0,1,2,\ldots).
\]

For all $n\geq0$ let $\omega_{n}(k)$ (cf. \href{https://oeis.org/A231599}{OEIS:A231599}
and \href{https://oeis.org/A333290}{A333290}) defined by
\[
\omega_{n}(k):=\begin{cases}
1 & \text{if }n,k=0,\\
0 & \text{if }n=0\text{ and }k>0,\\
\text{the coefficient of \ensuremath{x^{k}}}\text{ in the expansion of the product \ensuremath{{\displaystyle \prod_{j=1}^{n}(1-x^{j})}}} & \text{if }n>0\text{, }k=0,1,2,\ldots
\end{cases}
\]

For each $n\geq0$, we introduce the following sequences:

\begin{align*}
 & \mathbf{\omega}_{n}:=(\omega_{n}(0),\omega_{n}(1),\omega_{n}(2),\omega_{n}(3),\ldots), &  & \mathbf{p}:=(p(0),p(1),p(2),p(3),\ldots),\\
 & \mathbf{pic}:=(\text{pic}(0),\text{pic}(1),\text{pic}(2),\ldots), &  & \mathbf{p}_{{\scriptscriptstyle >}n}:=(p_{{\scriptscriptstyle >}n}(0),p_{{\scriptscriptstyle >}n}(1),p_{{\scriptscriptstyle >}n}(2),\ldots),\\
 & \mathbf{pic}_{n}:=(\text{pic}_{n}(0),\text{pic}_{n}(1),\text{pic}_{n}(2),\ldots), &  & \mathbf{p}_{{\scriptscriptstyle \leqq}n}:=(p_{{\scriptscriptstyle \leqq}n}(0),p_{{\scriptscriptstyle \leqq}n}(1),p_{{\scriptscriptstyle \leqq}n}(2),\ldots),
\end{align*}

where $p_{{\scriptscriptstyle \leqq}n}(k)$ denotes the number of
partitions of $k$ into parts $\leq n$ (or equivalently, into at
most $n$ parts).

($\mathbf{p}_{{\scriptscriptstyle >}n}$: cf. \href{https://oeis.org/A000041}{OEIS:A000041},
\href{https://oeis.org/A002865}{A002865}, \href{https://oeis.org/A008483}{A008483},
\href{https://oeis.org/A008484}{A008484}, \href{https://oeis.org/A185325}{A185325}-\href{https://oeis.org/A185329}{A185329})

($\mathbf{p}_{{\scriptscriptstyle \leqq}n}$: cf. \href{https://oeis.org/A000012}{OEIS:A000012},
\href{https://oeis.org/A004526}{A004526}, \href{https://oeis.org/A001399}{A001399},
\href{https://oeis.org/A001400}{A001400}, \href{https://oeis.org/A001401}{A001401},\href{https://oeis.org/A001402}{A001402},
\href{https://oeis.org/A026813}{A026813}-\href{https://oeis.org/A026816}{A026816})

By definition, the following limiting cases immediately follow from
the interpretations of these sequences: 
\begin{center}
$\omega_{0}=(1,0,0,0,\ldots)$, $\omega_{\infty}=\omega$, $\mathbf{p}=\mathbf{p}_{{\scriptscriptstyle >}0}=\mathbf{p}_{{\scriptscriptstyle \leqq}\infty}$,
$\mathbf{0}=(0,0,0,\ldots)=\mathbf{p}_{{\scriptscriptstyle >}\infty}=\mathbf{p}_{{\scriptscriptstyle \leqq}0}$,
$\mathbf{1}=(1,1,1,\ldots)=\mathbf{p}_{{\scriptscriptstyle \leqq}1}$.
\par\end{center}

We will use the well-known inverse relation

\[
\mathbf{\omega}_{n}^{-\ast}=\mathbf{p}_{{\scriptscriptstyle \leqq}n}\,\,(n=0,1,2,\ldots,\infty)
\]

where the superscript $\mathbf{a}^{-\ast}$ denotes the convolutional
inverse of a sequence $\mathbf{a}$.

Using these algebraic formulations and sequence definitions, the relations
established above for $\text{pic}_{n}(k)$ can be rewritten in a highly
compact, equivalent form via the following convolutions:
\begin{eqnarray*}
\mathbf{p}_{{\scriptscriptstyle >}n} & = & \mathbf{\omega}_{n}\ast\mathbf{p}\\
\mathbf{pic}_{n} & = & \mathbf{p}_{{\scriptscriptstyle >}n}\ast\mathbf{pic}
\end{eqnarray*}

To illustrate these algebraic properties, we explicitly compute the
initial terms of the sequences $\mathbf{pic}_{n}$ corresponding to
the cases $n=0,1,\text{and }2$.

\begin{align*}
\mathbf{pic}_{0} & =\mathbf{p}_{{\scriptscriptstyle >}0}\ast\mathbf{pic}=\mathbf{p}\ast\mathbf{pic}\\
 & =(1,1,2,3,5,7,11,15,22,30,42,56,77,101,\ldots)\ast(1,0,0,1,2,5,9,19,37,74,148,296,591,1183,\ldots)\\
 & =(1,1,2,4,8,16,32,64,128,256,512,1024,2048,4096,\ldots).
\end{align*}

\begin{align*}
\mathbf{pic}_{1} & =\mathbf{p}_{{\scriptscriptstyle >}1}\ast\mathbf{pic}\\
 & =(1,0,1,1,2,2,4,4,7,8,12,14,21,24,\ldots)\ast(1,0,0,1,2,5,9,19,37,74,148,296,591,1183,\ldots)\\
 & =(1,0,1,2,4,8,16,32,64,128,256,512,1024,2048,\ldots).
\end{align*}
\begin{align*}
\mathbf{pic}_{2} & =\mathbf{p}_{{\scriptscriptstyle >}2}\ast\mathbf{pic}\\
 & =(1,0,0,1,1,1,2,2,3,4,5,6,9,10,\ldots)\ast(1,0,0,1,2,5,9,19,37,74,148,296,591,1183,\ldots)\\
 & =(1,0,0,2,3,6,12,24,48,96,192,384,768,1536,\ldots).
\end{align*}

The following matrix summarizes the first $14$ terms of the sequences
$\mathbf{pic}_{n}$ computed for $n=0,1,\ldots,9$.

\[
\left[\begin{array}{cccccccccccccc}
1 & 1 & 2 & 4 & 8 & 16 & 32 & 64 & 128 & 256 & 512 & 1024 & 2048 & 4096\\
\noalign{\medskip}1 & 0 & 1 & 2 & 4 & 8 & 16 & 32 & 64 & 128 & 256 & 512 & 1024 & 2048\\
\noalign{\medskip}1 & 0 & 0 & 2 & 3 & 6 & 12 & 24 & 48 & 96 & 192 & 384 & 768 & 1536\\
\noalign{\medskip}1 & 0 & 0 & 1 & 3 & 6 & 10 & 21 & 42 & 84 & 168 & 336 & 672 & 1344\\
\noalign{\medskip}1 & 0 & 0 & 1 & 2 & 6 & 10 & 20 & 39 & 78 & 158 & 315 & 630 & 1260\\
\noalign{\medskip}1 & 0 & 0 & 1 & 2 & 5 & 10 & 20 & 38 & 76 & 152 & 305 & 610 & 1221\\
\noalign{\medskip}1 & 0 & 0 & 1 & 2 & 5 & 9 & 20 & 38 & 75 & 150 & 300 & 600 & 1201\\
\noalign{\medskip}1 & 0 & 0 & 1 & 2 & 5 & 9 & 19 & 38 & 75 & 149 & 298 & 595 & 1192\\
\noalign{\medskip}1 & 0 & 0 & 1 & 2 & 5 & 9 & 19 & 37 & 75 & 149 & 297 & 593 & 1187\\
\noalign{\medskip}1 & 0 & 0 & 1 & 2 & 5 & 9 & 19 & 37 & 74 & 149 & 297 & 592 & 1185\\
\noalign{\medskip}1 & 0 & 0 & 1 & 2 & 5 & 9 & 19 & 37 & 74 & 148 & 297 & 592 & 1184
\end{array}\right]
\]

By shifting each row of this matrix by its corresponding index $n=0,1,2,3,\ldots$,
we obtain precisely the matrix $\mathbf{R}$ discussed in Conjecture
1:

\[
\mathbf{R}=\left[\begin{array}{ccccccccccccccc}
1 & 1 & 2 & 4 & 8 & 16 & 32 & 64 & 128 & 256 & 512 & 1024 & 2048 & 4096 & \cdots\\
\noalign{\medskip}0 & 1 & 0 & 1 & 2 & 4 & 8 & 16 & 32 & 64 & 128 & 256 & 512 & 1024 & \cdots\\
\noalign{\medskip}0 & 0 & 1 & 0 & 0 & 2 & 3 & 6 & 12 & 24 & 48 & 96 & 192 & 384 & \cdots\\
\noalign{\medskip}0 & 0 & 0 & 1 & 0 & 0 & 1 & 3 & 6 & 10 & 21 & 42 & 84 & 168 & \cdots\\
\noalign{\medskip}0 & 0 & 0 & 0 & 1 & 0 & 0 & 1 & 2 & 6 & 10 & 20 & 39 & 78 & \cdots\\
\noalign{\medskip}0 & 0 & 0 & 0 & 0 & 1 & 0 & 0 & 1 & 2 & 5 & 10 & 20 & 38 & \cdots\\
\noalign{\medskip}0 & 0 & 0 & 0 & 0 & 0 & 1 & 0 & 0 & 1 & 2 & 5 & 9 & 20 & \cdots\\
\noalign{\medskip}0 & 0 & 0 & 0 & 0 & 0 & 0 & 1 & 0 & 0 & 1 & 2 & 5 & 9 & \cdots\\
\noalign{\medskip}0 & 0 & 0 & 0 & 0 & 0 & 0 & 0 & 1 & 0 & 0 & 1 & 2 & 5 & \cdots\\
\noalign{\medskip}0 & 0 & 0 & 0 & 0 & 0 & 0 & 0 & 0 & 1 & 0 & 0 & 1 & 2 & \cdots\\
\noalign{\medskip}0 & 0 & 0 & 0 & 0 & 0 & 0 & 0 & 0 & 0 & 1 & 0 & 0 & 1 & \cdots
\end{array}\right]
\]

By combining the two fundamental convolution formulas above, we immediately
obtain the following identity for $\mathbf{pic}_{n}$:

\[
\mathbf{pic}_{n}=\mathbf{p}_{{\scriptscriptstyle >}n}\ast\mathbf{pic}=\mathbf{\omega}_{n}\ast\left(\mathbf{p}\ast\mathbf{pic}\right)=\mathbf{\omega}_{n}\ast(1,1,2,4,8,16,32,\ldots)
\]

Convolving both sides of this equation with $\mathbf{\omega}_{n}^{-\ast}$,
we arrive at:

\[
\mathbf{\omega}_{n}^{-\ast}\ast\mathbf{pic}_{n}=\mathbf{p}_{{\scriptscriptstyle \leqq}n}\ast\mathbf{pic}_{n}=\mathbf{p}\ast\mathbf{pic}=(1,1,2,4,8,16,32,\ldots)
\]

Alternatively, convolving both sides with $\mathbf{p}^{-\ast}=\mathbf{\omega}$
we get the remarkably symmetric relation:

\[
\mathbf{\omega}\ast\mathbf{pic}_{n}=\mathbf{\omega}_{n}\ast\mathbf{pic}
\]


\section{The Pentagonal Transform and Its Recurrence Matrix}

Let $S$ denote the left shift operator acting on sequences, defined
by 
\[
S((a_{n})_{n\ge0})=(a_{n+1})_{n\ge0}.
\]

For each $k\ge0$, let $\mathbf{e}_{k}=(e_{k}(n))_{n\ge0}$ denote
the $k$-th standard basis sequence, defined via the Kronecker delta
by $e_{k}(n):=\delta_{k,n}\quad(n\ge0).$ In particular, $\mathbf{e}_{0}=(1,0,0,\dots)$
and $\mathbf{e}_{1}=(0,1,0,\dots)$, where the only non-zero term
of $\mathbf{e}_{k}$ appears at index $k$. We will frequently use
throughout this chapter the sequence $\mathbf{\delta}=(\delta(n))_{n\ge0}$
\[
\mathbf{\delta}:=S(\mathbf{e}_{0}-\mathbf{\omega})=(1,1,0,0,-1,0,-1,0,0,0,0,1,0,0,1,0,\dots).
\]

That is, $\delta(n)=$\href{https://oeis.org/A257628}{A257628(n+1)}
$(n\geq0)$. We retain the notation for the vectors $\mathbf{\omega}_{n}$,
$\mathbf{p}$, $\mathbf{p}_{{\scriptscriptstyle \leqq}n}$, $\mathbf{p}_{{\scriptscriptstyle >}n}$,
$\mathbf{pic}$, $\mathbf{pic}_{n}$ established in the previous chapter.
\begin{defn*}
Let $\mathbf{a}=(a(n))_{n\ge0}$ be a fixed sequence with $a(0)=1$.
We associate with $\mathbf{a}$ an infinite two-dimensional matrix
$\mathbf{R}(\mathbf{a})=(r(i,j))$ ($i,j\in\mathbb{Z}$) by the following
recurrence relation:
\[
r(i,j)=\begin{cases}
0 & \text{if }i<0\text{ or }j<i,\\
a(j) & \text{if }i=0\text{ and }j\ge0,\\
r(i-1,j-1)-r(i-1,j-1-i) & \text{otherwise.}
\end{cases}
\]
The \emph{pentagonal transform} of the sequence $\mathbf{a}$, denoted
by $\mathcal{P}(\mathbf{a})$ is the sequence $\mathbf{b}=(b(n))_{n\ge0}$
whose $n$-th term is given by the sum of the elements in the $n$-th
column of $\mathbf{R}(\mathbf{a})$:
\[
b(n):=\sum_{i=0}^{n}r(i,n)\quad(n\ge0).
\]
\end{defn*}
It is an immediate consequence of the definition that the transform
$\mathcal{P}$ is linear. The pentagonal transform $\mathcal{P}$
and its underlying recurrence relation possess widespread applications;
in fact, the wealth of combinatorial identities derivable from them
could fill an entire volume on their own. For our present purposes,
however, it suffices to note that for any arbitrary sequence $\mathbf{a}$,
we have
\[
\mathcal{P}(\mathbf{a})=\mathbf{\delta}*\mathbf{a},
\]
 where $*$ denotes the convolution of sequences. 

To illustrate the mechanics of the pentagonal transform $\mathcal{P}$
and the structure of its associated matrix $\mathbf{R}(\mathbf{a})$,
we present two concrete examples generated via a custom JavaScript
application. Since these matrices are inherently infinite, the figures
display only their first few rows and columns. The first displays
matrix $\mathbf{R}(\mathbf{e}_{0})$ corresponding to the standard
basis sequence $\mathbf{e}_{0}=(1,0,0,\dots)$ (indicated in the top
blue row). As expected, the resulting column sums (highlighted in
the bottom green row) precisely yield the sequence $\mathbf{\delta}$,
reflecting the property $\mathcal{P}(\mathbf{e}_{0})=\mathbf{\delta}$.
Similarly, the second demonstrates the transform applied to the sequence
of partition numbers $\mathbf{p}=(p(n))_{n\ge0}$. In this case, the
top blue row contains the initial partition values, while the bottom
green row displays the row sums constituting the transformed sequence
$\mathcal{P}(\mathbf{p})=S(\mathbf{p})$. 
\begin{center}
\includegraphics[width=0.42\paperwidth]{rek_1} \hspace{2mm}\includegraphics[width=0.42\paperwidth]{rek_p}
\par\end{center}

Furthermore, the next figure illustrates the transform applied to
the sequence $\mathbf{a}=\mathbf{p}*\mathbf{pic}=(1,1,2,4,8,16,\dots)$.
This particular case is of primary interest as it is directly related
to Conjecture 1, by virtue of the fact that each value of the sequence
$\mathcal{P}(\mathbf{a})$ displayed in the bottom green row is precisely
one greater than the corresponding value of $\mathcal{L}(n)$.

\hfill{}\includegraphics[width=0.86\textwidth]{rek_2}\hfill{}

If $\mathbf{1}$ denotes the constant sequence $(1,1,1,\dots)$ and
$\mathbf{L}$ denotes the sequence $(\mathcal{L}(2),\mathcal{L}(3),\mathcal{L}(4),\dots)$,
then this relationship can be expressed in the following compact form
\[
\mathbf{L}=\mathcal{P}(\mathbf{a})-\mathbf{1}=\delta*\mathbf{p}*\mathbf{pic}-\mathbf{1},
\]
 from which the formula stated in Conjecture 1 can be obtained via
basic algebraic manipulations. 

However, the connection runs much deeper; it manifests not only at
the level of the transformed sequence $\mathbf{L}=\mathcal{P}(\mathbf{a})-\mathbf{1}$,
but also throughout the entire recurrence matrix, appearing in every
single row. Consequently, this observation opens a clear avenue toward
proving Conjecture 1. To elucidate this deeper connection, we introduce
a so-called rank-increment matrix associated with the derivation relation,
and compare its empirically observed values with the matrix of the
recently discussed transform. In alignment with the notation established
in Chapter 3, let $\mathcal{A}_{n}:=x\mathfrak{H}_{n-2}y$ denote
the set of admissible words of length $n\geq2$, and let $\partial_{m}(\mathcal{A}_{n}):=\{\partial_{m}(w):w\in\mathcal{A}_{n}\}$
for any $m\geq1$. For notational convenience, we index the rows of
the rank-increment matrix $\mathbf{R}=(r_{m,n})$ starting from $1$
and its columns starting from $2$, so that $m\ge1$ and $n\ge2$.
The entries $r_{n,m}$ of the rank-increment matrix are formally defined
as follows: 
\begin{enumerate}
\item The entry in the first row and $n$-th column is given by the rank
of the vector system $\partial_{1}(\mathcal{A}_{n-1})$; that is,
\[
r_{1,n}=\text{rank}\left(\partial_{1}(\mathcal{A}_{n-1})\right)\quad\text{for all }n\ge2.
\]
\item For the subsequent rows where $m\ge2$, the entry $r_{m,n}$ is defined
as the rank increment obtained by adjoining the vector system $\partial_{m}(\mathcal{A}_{n-m})$
to the combined system 
\[
\bigcup_{k=1}^{m-1}\partial_{k}(\mathcal{A}_{n-k}).
\]
 That is,
\[
r_{m,n}=\text{rank}\left(\bigcup_{k=1}^{m}\partial_{k}(\mathcal{A}_{n-k})\right)-\text{rank}\left(\bigcup_{k=1}^{m-1}\partial_{k}(\mathcal{A}_{n-k})\right)
\]
 for all $m\ge2$ and $n\ge m+1$. 
\end{enumerate}
This formulation is completely equivalent to the rank-increment procedure
discussed in Chapter 3; however, while the focus here is on conceptual
clarity and conciseness, the description there was tailored to precisely
reflect the algorithmic implementation. By exploiting the underlying
duality, for rows where $2\leq m\leq n-3$, we restricted our analysis
to the vector systems $\partial_{m}(x\mathcal{A}_{n-m-1})$, which
comprise only half of the full system $\partial_{m}(\mathcal{A}_{n-m})$.
Furthermore, those elements were adjoined to the matrix sequentially,
one by one, according to a fixed ordering, discarding those that did
not increase the rank. 

The empirically computed rank-increment matrix is presented below,
exhibiting a perfect match with the strictly upper triangular part
of the transformation matrix $\mathbf{R}(\mathbf{a})$ initialized
with the sequence $\mathbf{a}=(1,1,2,4,8,16,\dots)$. 

\[
\begin{array}{c|cccccccccccccccc}
{\color{lightgray}} & 2 & 3 & 4 & 5 & 6 & 7 & 8 & 9 & 10 & 11 & 12 & 13 & 14 & 15 & 16 & 17\\
\hline \partial_{1} & 1 & 1 & 2 & 4 & 8 & 16 & 32 & 64 & 128 & 256 & 512 & 1024 & 2048 & 4096 & 8192 & 16384\\
\partial_{2} &  &  & 0 & 1 & 2 & 4 & 8 & 16 & 32 & 64 & 128 & 256 & 512 & 1024 & 2048 & 4096\\
\partial_{3} &  &  &  & 0 & 0 & 2 & 3 & 6 & 12 & 24 & 48 & 96 & 192 & 384 & 768 & 1536\\
\partial_{4} &  &  &  &  & 0 & 0 & 1 & 3 & 6 & 10 & 21 & 42 & 84 & 168 & 336 & 672\\
\partial_{5} &  &  &  &  &  & 0 & 0 & 1 & 2 & 6 & 10 & 20 & 39 & 78 & 158 & 315\\
\partial_{6} &  &  &  &  &  &  & 0 & 0 & 1 & 2 & 5 & 10 & 20 & 38 & 76 & 152\\
\partial_{7} &  &  &  &  &  &  &  & 0 & 0 & 1 & 2 & 5 & 9 & 20 & 38 & 75\\
\partial_{8} &  &  &  &  &  &  &  &  & 0 & 0 & 1 & 2 & 5 & 9 & 19 & 38\\
\partial_{9} &  &  &  &  &  &  &  &  &  & 0 & 0 & 1 & 2 & 5 & 9 & 19\\
\partial_{10} &  &  &  &  &  &  &  &  &  &  & 0 & 0 & 1 & 2 & 5 & 9\\
\partial_{11} &  &  &  &  &  &  &  &  &  &  &  & 0 & 0 & 1 & 2 & 5\\
\partial_{12} &  &  &  &  &  &  &  &  &  &  &  &  & 0 & 0 & 1 & 2\\
\partial_{13} &  &  &  &  &  &  &  &  &  &  &  &  &  & 0 & 0 & 1\\
\partial_{14} &  &  &  &  &  &  &  &  &  &  &  &  &  &  & 0 & 0\\
\partial_{15} &  &  &  &  &  &  &  &  &  &  &  &  &  &  &  & 0\\
\partial_{16}
\end{array}
\]

The fact that the powers of two appear in the first row of the empirical
rank-increment matrix $\mathbf{R}$ is theoretically straightforward
to verify; specifically, $\text{rank}(\partial_{1}(\mathcal{A}_{n}))=2^{n-3}$
for $n\ge3$. Indeed, the cardinality of the underlying set is given
by $|\mathcal{A}_{n-1}|=|x\mathfrak{H}_{n-3}y|=2^{n-3}$. To establish
the linear independence of the vector system $\partial_{1}(\mathcal{A}_{n-1})$,
it suffices to examine the first term of the derivative of each word
$w\in\mathcal{A}_{n}$ with respect to the ordering fixed in Chapter
3. Consequently, to complete the proof of Conjecture 1, it suffices
to demonstrate that under the appropriate reindexing, the defined
rank increments $r_{m,n}$ satisfy the recurrence relation prescribed
in the definition of the pentagonal transform. 

To further investigate the underlying algebraic foundations, it is
highly instructive that the matrix can be viewed as a two-dimensional
extension $\mathbf{pic}_{n}(k)$ of the sequence of pure inverting
compositions $\mathbf{pic}(k)$, as discussed in detail in the preceding
section. In view of what is described there, \emph{Conjecture 1 would
follow from directly showing that} the vector system $\partial_{m}(\mathcal{A}_{k})$
increases the rank of the system 
\[
\bigcup_{j=1}^{m-1}\partial_{m-j}(\mathcal{A}_{k+j})=\partial_{m-1}(\mathcal{A}_{k+1})\cup\partial_{m-2}(\mathcal{A}_{k+2})\cup\dots\cup\partial_{1}(\mathcal{A}_{k+m-1})
\]
 by exactly 
\[
\sum_{i=0}^{k}p_{{\scriptscriptstyle >}m}(i)\cdot\text{pic}(k-i)=p_{{\scriptscriptstyle >}m}(0)\,\text{pic}(k)+p_{{\scriptscriptstyle >}m}(1)\,\text{pic}(k-1)+\cdots+p_{{\scriptscriptstyle >}m}(k-1)\,\text{pic}(1)+p_{{\scriptscriptstyle >}m}(k)\,\text{pic}(0).
\]


\section{Other Conjectures on Derivation Relations}

Recall that for any admissible word $w\in\mathfrak{H}_{0}$ and any
integers $a,b\geq0$, the shuffle regularization of the word $y^{a}wx^{b}$
is given by
\[
\text{reg}_{\ensuremath{\shuffle}}(y^{a}wx^{b}):=\sum_{i=0}^{a}\sum_{j=0}^{b}(-1)^{i+j}\,y^{i}\shuffle\left(y^{a-i}wx^{b-j}\right)\shuffle x^{j}
\]

\begin{conjecture}
(Generalized Derivation Relation) For any admissible word $w\in\mathfrak{H}_{0}$
and any integers $n>0$, $a,b\geq0$, and $m>a+b$, we have
\[
\zeta\left(\text{reg}_{\ensuremath{\shuffle}}\left[\partial_{n}^{m}\left(y^{a}wx^{b}\right)\right]\right)=0.
\]
\end{conjecture}

Moreover, the quasi-derivation relation can be generalized in a similar
manner:
\begin{conjecture}
(Generalized Quasi-derivation Relation) For any rational number $c\in\mathbb{Q}$
, any admissible word $w\in\mathfrak{H}_{0}$ and any integers $n>0$,
$a,b\geq0$, and $m>a+b$, we have
\[
\zeta\left(\text{reg}_{\ensuremath{\shuffle}}\left[\left(\partial_{n}^{(c)}\right)^{m}\left(y^{a}wx^{b}\right)\right]\right)=0.
\]
\end{conjecture}

The bold rows of the table (which also incorporate rows from the table
in \cite{TANAKA20092021}) below present the number of linearly independent
relations obtained from Conjecture 4, Conjecture 5 (calculated via
a JavaScript application developed by the author), and the linear
part of Kawashima's relation, respectively. 

\parbox[t]{1\columnwidth}{%
\begin{center}
Table 5. Number of Independent Relations for MZVs
\par\end{center}
\begin{center}
\begin{tabular}{|c|c|c|c|c|c|c|c|c|c|c|c|c|c|}
\hline 
{\small{}$\vphantom{\dfrac{A}{B}}n$} & {\small{}3} & {\small{}4} & {\small{}5} & {\small{}6} & {\small{}7} & {\small{}8} & {\small{}9} & {\small{}10} & {\small{}11} & {\small{}12} & {\small{}13} & {\small{}14} & {\small{}15}\tabularnewline
\hline 
\hline 
{\small{}$\vphantom{\dfrac{A}{B}}\partial_{n}$} & {\small{}1} & {\small{}2} & {\small{}5} & {\small{}10} & {\small{}22} & {\small{}44} & {\small{}90} & {\small{}181} & {\small{}363} & {\small{}727} & {\small{}1456} & {\small{}2912} & {\small{}5825}\tabularnewline
\hline 
{\small{}$\vphantom{\dfrac{A}{B}}\partial_{n}^{m}$} & \textbf{\small{}1} & \textbf{\small{}2} & \textbf{\small{}5} & \textbf{\small{}10} & \textbf{\small{}24} & \textbf{50} & \textbf{104} & \textbf{213} & \textbf{431} & \textbf{873} & \textbf{1756} & \textbf{3532} & \textbf{7080}\tabularnewline
\hline 
\hline 
{\small{}$\vphantom{\dfrac{A}{B}}\partial_{n}^{(c)}$} & 1 & 2 & 5 & 10 & 23 & 46 & 98 & 200 & 410 & 830 & 1679 & - & -\tabularnewline
\hline 
{\small{}$\vphantom{\dfrac{A}{B}}\left(\partial_{n}^{(c)}\right)^{m}$} & \textbf{1} & \textbf{2} & \textbf{5} & \textbf{11} & \textbf{24} & \textbf{50} & \textbf{105} & \textbf{217} & \textbf{445} & \textbf{910} & - & - & -\tabularnewline
\hline 
\hline 
$\vphantom{\dfrac{A}{B}}$Ohno & 1 & 2 & 5 & 10 & 23 & 46 & 98 & 199 & 411 & 830 & 1691 & - & -\tabularnewline
\hline 
$\vphantom{\dfrac{A}{B}}$Lin.Kaw. & \textbf{1} & \textbf{2} & \textbf{5} & \textbf{10} & \textbf{23} & \textbf{46} & \textbf{98} & \textbf{200} & \textbf{413} & \textbf{838} & \textbf{1713} & - & -\tabularnewline
\hline 
$\vphantom{\dfrac{A}{B}}$All & 1 & 3 & 6 & 14 & 29 & 60 & 123 & 249 & 503 & 1012 & 2032 & 4075 & 8164\tabularnewline
\hline 
\end{tabular}
\par\end{center}%
}

\bibliographystyle{plain}
\addcontentsline{toc}{section}{\refname}\nocite{*}
\bibliography{conjecture}

\end{document}
